\documentclass[conference]{IEEEtran}
\IEEEoverridecommandlockouts
% The preceding line is only needed to identify funding in the first footnote. If that is unneeded, please comment it out.
%Template version as of 6/27/2024

\usepackage{acronyms}
\usepackage{commands}

% \usepackage{cite}
\usepackage{amsmath,amssymb,amsfonts}
\usepackage{algorithmic}
\usepackage{graphicx}
\usepackage{textcomp}
\usepackage{xcolor}
\def\BibTeX{{\rm B\kern-.05em{\sc i\kern-.025em b}\kern-.08em
    T\kern-.1667em\lower.7ex\hbox{E}\kern-.125emX}}

\usepackage[backend=biber,style=ieee]{biblatex}
\addbibresource{HPC_seminar.bib}

\begin{document}

\title{A Cost-Benefit Study of Template Metaprogramming}

\author{\IEEEauthorblockN{1\textsuperscript{st} Miguel Veganzones Parellada}
\IEEEauthorblockA{\textit{Technical University of Munich} \\
% \textit{Technical University of Munich}\\
Munich, Germany \\
miguel.veganzones@tum.de}
}

\maketitle

\begin{abstract}
\end{abstract}

\begin{IEEEkeywords}
component, formatting, style, styling, insert.
\end{IEEEkeywords}

\section{Introduction}
explore how increasing compiler visibility
through templates and template metaprogramming techniques affects the
performance and generated code quality of numerical kernels.

will compare implementations of the same algorithm with various degrees of
runtime/compile time information, focusing on runtime, compile time, code size,
and constant propagation in the generated assembly.

The study will compare runtime-generic implementations against progressively
more compile-time-specialized variants.

C++ philosophy:
The aim is to allow a programmer to work at the highest feasible level of ab-
straction by providing
- A simple and direct mapping to hardware
- Zero-overhead abstraction mechanisms

Various forms of user-specified compile-time computation are essential in
critical embedded systems applications, much low-level code, and many high-
end numerical applications.
\cite[p.~6]{stroustrup_foundations_2011}

\section{Compile-Time Information}

Template metaprogramming allows implementing thin layers of
abstraction whose complexity is resolved at compile time.

The extent to which this can occur is limited by compiler visibility and
available information.

Compiler visibility is heavily influenced by dynamic binding.
Compile time information must be hard-coded into the source, and it must also be
propagated at compile time to be usable in \texttt{constexpr} contexts.
Compile time information can be lowered to runtime information in a nested
context, reducing the compiler's visibility in said context.
This process is irreversible, and defines a flow of information from its source
to its destination, where it is eventually consumed.
External inputs are never present in the source and are therefore never usable
at compile time.
Only information present in the source, and computable results can be used at
compile time.

Template metaprogramming allows manipulating values and types at compile time.
Values in constexpr contexts additionally.
Types are very powerful in a statically, strongly typed language

This defines a unique direction of the flow of information from its source to
its sink, where it is eventually consumed.

If a value is ever lowered to the runtime, it cannot be elevated to the runtime
at a later stage as the information is lost from the present context and all
subsequent contexts.

Template metaprogramming is a powerful modern \cpp{} tool, which is leveraged to
afford efficient abstractions in highly-optimized applications \cite{hendricks_lessons_2025}.

The effects of compiler visibility and template metaprogramming are often
understood in terms of zero-cost abstractions, but the implication go beyond
that, and are describe below.

Simple and direct mapping to hardware and zero-cont abstractions
\cite[p.~2]{stroustrup_foundations_2011}


Comnpute a type at compile time
\cite[p.~6]{stroustrup_foundations_2011}


\subsection{Zero-Cost Abstractions}
C++ follows the Zero-Overhead Principle
The zero-overhead principle is a C++ design principle that states:

    You don't pay for what you don't use.
    What you do use is just as efficient as what you could reasonably write by hand.

In general, this means that no feature should be added to C++ that would impose any overhead, whether in time or space, greater than a programmer would introduce without using the feature.

The only two features in the language that do not follow the zero-overhead principle are runtime type identification and exceptions, and are why most compilers include a switch to turn them off.
Zero-cost abstractions is a core design philosophy of the \cpp{} language

\subsection{Compile Time Work}

4 Compile-time Computation
Sometimes, we prefer a computation be done at compile-time. The reasons
vary, for example:
• Efficiency: To pre-calculate a value (often a size). For simple cases, that is
done by an optimizer. Examples include object and array sizes and table
values.
• Type-safety: To compute a type at compile time.
• Simplify concurrency: you can’t have a race condition on a constant.

\cite[p.~6]{stroustrup_foundations_2011}

Templates as Type function: template functions generate functions, template classes
generate classes
\cite[p.~21]{stroustrup_foundations_2011}

\subsubsection{Move computation from runtime to compile time}
Moving any computation from runtime to compile time should increase compile
times in favor of runtime.
Wether to do it comes down to whether it is possible and what cost is more
important.
Compile time execuition is slower as it executes on an interpreter inside the
compiler.

\subsubsection{Avoid redundant computation at runtime}
(some of the computation may be done repeatedly)
API designers, for ease of usability may desire to repeat computation to
simplify the callsite
In our case later on, an efficient implemenation may have:
f(in A, in B, out C) with a pre-constructed C
instead of 
C = f(in A, in B)
As calculating the structure of C would need to be done per call, rather than
per caller.

\subsubsection{Compile-Time Complexity}
Cost of recompilation/code size vs runtime

reduced over time
from explicit sfinae to 
if constexpr
template for

\subsection{Code Generation}

\subsubsection{Enable optimizations}
    - automatic unrolling
    - loop stamping

\subsubsection{Code Size}
    - compile time logic to compute values is not preserved in the executable
Reduce object sizes
    - reduced indirections and stuff
    Hard coding a result may be cheaper than hard coding the procedure to obtain
    it
Increase code size
    Template deduplication
    Inlining/unrolling

\subsubsection{Type simplification/stronger properties/weaker properties/requirements}
    - Trivial types
    - Standard layout
    - Trivial copyability

\subsubsection{Weaker move}
Pointer/reference semantics vs value semantics

\subsubsection{constexpr enabling and propagation}
Initialization

\subsubsection{Compile time contracts}

\subsection{Architectural}

\subsubsection{Reduce memory operations}
Immediate arguments are inevitably fetched with their instruction, bypassing
loads from the data pipeline to be fetched.
This reduces the overall number of load operations in the assembly.
What type categories and bit widths are allowed is architecture-dependent.
intel64 accepts small integral values and no floating point, while amd64 also
supports floatng point immediates to some extent I believe

\subsubsection{Reduce function arguments}
Static global data region has predictable locations
Heap is clearly in a disadvantage
Stack arguments must be input arguments as their location is unknown to an
arbitrary caller.
Static/global symbols in \texttt{.rdata} or \texttt{.data} have predictable
locations, which can potentially be hardcoded in the load instruction.

\subsubsection{Reduce register pressure}
Immediate arguments reduce register pressure as they never become part of the
in-flight or architectural state, but are rather consumed and discarded.

\subsubsection{Read-only data and coherence protocols}
Read only data regions are simpler to manage by a multithreaded architecture as
it promises stronger invariants.
Race conditions, initlaization

\subsection{Memory Regions}

\subsubsection{Reduce the need for allocation}
The stack is compiler controlled.
It is restricted to objects of known size and local-context-determined lifetime.
The heap softens these restrictions in exchange for \ac{os} interaction,
increased compiler uncertainty and indirections.
Lowering values from runtime to compile time may enable the use of more
restrictive sets of memory

Heap: Free
Stack: Known structure, dynamically evaluated
Static data: Statically known layout (v2 sizes)
Constexpr data: Statically known values (v3 sizes)


Only if we ex-
plicitly allocate an (unnamed) Point on the free store (the heap), as done for
the Point pointed to by p, do we incur memory overhead (and allocation over-
head). Such very simple user-defined types are critical to type-rich program-
ming and very common
\cite[p.~3]{stroustrup_foundations_2011}

\subsubsection{Static/Global Data}
Increase data size
    Static constexpr

\subsubsection{Increased stack pressure}
possible oveflow

\subsubsection{Reduce indirections from heap to the stack}

Increase compile time complexity
    If scaling is bad, the compile time interpreter can exhaust resources

\section{Design Space}
runtime vs compile-time parameters
flexibility vs performance
abstraction vs generated code simplicity

\section{Compute Kernel}

The extent to which compile time information can be exploited effectively is
highly application-dependent.
In particular, a reference application must be affected in a meaningful way by
the absence or presence of this compile time information.
Tensor contraction is the algorithm chosen as reference application as there are
two clear progression steps that can be exploited both in both the container and
algorithm implementations.

\cite{hendricks_lessons_2025}

The microkernels

Compile-time information is and the reference algorithm must be able to
effectively exploit it

numerical kernels as interesing applications

Types - would be pathological to consider it
Rank
Sizes

Tensor contraction 

qualitative and quantitative changes

dynamnic to compile time polymorphism

memory allocation
immediate variables


v1: not knowing the rank effectively requires a second loop during iteration,
effectively a multi-index

\section{Implementation}

We consider three levels of compile time information that the implementation can
exploit in a meaningful way.
Each level will build on top of the previous one, always elevating information
from runtime to compile time.
These levels are denominated \textit{runtime rank}, \textit{compile-time size}
and \textit{compile-time sizes} (\level{v1}, \level{v2}, \level{v3}
respectively), according to how the available information affects each
implementation.

These can be though of as different specification standards with different
requirements for compliance.

\subsection{Runtime Rank (\level{v1})}

This level specifies the lowest degree of compile-time information expected by
each of the implementations.
It requires static typing and static binding.
This is a sensible baseline as it is often taken for granted.

Static typing implies that the concrete types being used are known to the
compiler (\ie{} no \texttt{std::any} or \texttt{void*}).
Static binding implies that function calls are resolved at compile time (\ie{}
no virtual dispatch or function pointers).

In this context, the compiler has full visibility of the concrete
types and call structure, enabling optimization across function boundaries and
efficient code generation.
However, no algorithm-specific information such as ranks, dimensions or
contraction patterns will be accessible during code generation.

\subsubsection{\texttt{v1::tensor} (\level{t1})}

This variant must use the heap to store both its data and its metadata
The data buffer will be a contiguous buffer of \texttt{T}.
Similarly, the metadata buffer will contain the tensor sizes and strides for
efficient access with a multi-index.
Accessing these buffers through a pointer indirection does not meaningfully
affect runtime on its own, as compared to the heap.
However, we will have 7 different allocations\footnote{Two allocations for each
    tensor and one for the contraction index set. Each tensor uses two buffers
    in the heap, these could be merged into one, but the additional complexity
was deemed unimportant.} which will demand more memory subsystem resources.

\subsubsection{\texttt{v1::contraction} (\level{c1})}

Since the rank of the tensors is unknown at compile time, the first-level
contraction implementation uses a multi-index to iterate over the tensors.
The multi-index uses an internal \texttt{for} loop for each increment,
introducing significant overhead in iteration.

\subsection{Compile-Time Rank (v2)}

The second level requires that the rank of the tensors and the contraction
degree be known at compile time.

\subsubsection{\texttt{v2::tensor} (\level{t2})}

Knowing the rank of the tensor allows the strides and sizes of the tensors
involved to be stored within the tensor object itself.
Thanks to this compile time information, the metadata of the tensor will be
stored in the stack, rather than the heap, and accessible without a pointer
indirection.\footnote{Moving data from the stack to the heap also removes a
costly memory allocation, but these are not included in the benchmark loop.}
The data itself must still be heap-allocated, however, as the size of each
dimension is not known.

\subsubsection{\texttt{v2::contraction} (\level{c2})}

Knowing the rank of the tensors significantly simplifies iterating over the
tensors\footnote{The rank of a tensor corresponds to the number of nested for
loops required to iterate over it in a structured way.}.
This implementation sets up two sets of nested \texttt{for} loops using template
metaprogramming, removing the need to iterate using a multi-index.

\subsection{Compile-Time Sizes (v3)}

The third level requires that the sizes of the tensors and the indices to
contract over, be known at compile time.

\subsubsection{\texttt{v3::tensor} (\level{t3})}

This affects the tensor implementation in two major ways.
For once, the data in the tensor may be stored in the stack, significantly
improving locality, specially for small tensors.
Additionally, tensor sizes and strides become part of the type, reducing the
size of each instance.

\subsubsection{\texttt{v3::contraction} (\level{c3})}

Knowing the size of the tensors allows propagating these values into the loop
limits into the two sets of nested \texttt{for} loops.
This enables inlining the loop limits into the assembly as integral immediate
arguments, and enables aggressive loop unrolling and vectorization.

Additionally, iteration functions now require less input arguments, since the
location of static data is known at compile time.
This change makes feasible computing the target linear index into each tensor
during iteration, removing an inner product computation between the strides and
the multi-index on each tensor access.\footnote{This same optimization was
unsuccessful for the \texttt{v2::contraction} iteration scheme, presumably for
the significant increase in function arguments.}

\subsection{Signatures}
The first level tensor has one template parameter for its value type
\begin{itemize}
    \item \texttt{v1::tensor<T>}
    \item \texttt{v2::tensor<T, N>}
    \item \texttt{v3::tensor<T, A1, ..., AN>}
\end{itemize}
\begin{itemize}
    \item \texttt{v1::contract<T>(in A, in B, out C, in cis)}
    \item \texttt{v2::contract<T, N>(in A, in B, out C, in cis)}
    \item \texttt{v3::contract<T, A1, ..., AN, B1, ... BN, C1, ... CN, CIS>(in A, in B, out C)}
\end{itemize}

\subsection{Compatibility}

Contraction and tensor implementations of different levels are compatible
because each level extends the previous one with stricter requirements and
compile time information can be lowered to runtime information.
Since the opposite is not true, compatibility goes in one direction, making
\texttt{v3::tensor} compatible with any contraction implementation, 
\texttt{v2::tensor} compatible with \texttt{v2::contraction} and
\texttt{v1::contraction}, and \texttt{v1::tensor} only compatible with
\texttt{v1::contraction}.

This will be used to assess what part of the performance improvements can be
attributed to changes in the tensor implementation or the contraction
implementation.

% t3 meets the requirements, but not the limitations of v1, and can, therefore, be used by c1.

\section{Experimental Setup}

arch linux 6.18.35-1-lts
GCC 16.1.1 c++26 mode
i7 9750h
32kb L1
256 KB L2
12 MB L3 unified

\section{Results}
Introduce metrics
\subsection{Runtime}
\subsection{Compile Time}
\subsection{Code Generation}
\section{Discussion}
\section{Conclusion}



\printbibliography

\end{document}
